\section{Introduction}

A long-horizon agent has to decide which past experiences to retain, which to
retrieve for a task, and which to expose to the model that answers it. For
retention, the relevant quantity is a memory's expected effect over future
tasks; for exposure, it is the effect on the task at hand. Both depend on how a
memory's effects vary across tasks: whether it helps many tasks or few, and
whether a memory that helps one task can hurt another. Memory systems already score items~\cite{multifactor,memoryrone,umem}, and prior
work observes that value depends on the task~\cite{ember,evomemeval,meclear}. We provide a measurement protocol that jointly holds the filler context
fixed, keeps the writer blind to the later task, and records the memory's
cost beside its effect. We apply it to web-agent memory on
LongMemEval-V2~\cite{lmev2}. The central finding is that usefulness depends
on the memory--task pair and on whether the representation preserves the
content needed for that task.

Consider a hypothetical past page visit that shows a project's budget. A
task-blind writer stores that visit as a memory without seeing any future
question. A later question asks for the budget, and the same reader answers it
twice, with the same filler, once with and once without the stored memory; the
reward difference is that memory's effect on that question. Repeating the
comparison for other questions and representations yields the memory's
task-conditioned profile.

The controlled comparisons show that a memory gains more on tasks for
which it contains answer-bearing information than on matched controls.
The tested summaries and structured records gain less than Raw, including
when total context is fixed. These findings make preserved task utility a concrete evaluation target
for memory compression, alongside its costs: a shorter memory is useful
only if it still helps the reader answer the relevant question. They also expose two
practical challenges: generic retrieval misses most screen-defined
answer-bearing memories, and individual memory--task labels can change
across measurement conditions. The aggregate bearing--control contrast
remains positive across the tested sessions, positions and readers;
that stability does not make each individual label stable.

\paragraph{Scope.}
Our empirical evidence comes from single-call factual-recall evaluation: the
memories are page visits from web-agent trajectories, the future evaluation
tasks are questions with deterministically scored answers, and a reader answers
each question in one call without acting in an environment. We make no claim
about procedural memory, acting agents, or a trained retention policy.

\paragraph{Contributions.}
\textbf{(1) A conditional memory utility protocol} (Section~\ref{sec:method}):
paired evaluation with identical context filler, a memory writer that never
observes the future task, and a cost ledger that records formation, storage,
retrieval and exposure costs beside each measurement.
\textbf{(2) Bounded findings for factual recall} (Sections~\ref{sec:screened},
\ref{sec:similarity}, \ref{sec:repr}; Figures~\ref{fig:robust}
and~\ref{fig:mech}). Within the same memory, answer-bearing task pairs gain
$\lfHeadD$ $[\lfHeadLo, \lfHeadHi]$ more reward than matched-control pairs. The
evaluated task-blind compressed representations show smaller gains than Raw at
their own length and at a fixed total context, consistent with reduced
retention of answer-bearing text, and generic retrieval recovers only part of
the screen-defined answer-bearing set ($\soRChkTenPct\%$ in its top ten, against
about $\soChanceTen\%$ by chance).
\textbf{(3) Measurement-reliability findings} (Sections~\ref{sec:metrics},
\ref{sec:readertransfer}, \ref{sec:r0}, \ref{sec:limitations}). The headline
contrast keeps its direction across a new session with randomized position
($\soXoneD$) and a second reader ($\rd{\rgBmcRawGem}$), while its magnitude and
per-pair labels change; sign-flip and harm summaries are confounded by pairs per
memory and baseline correctness, and we show how to condition both away.
\textbf{(4) A dataset} (Section~\ref{sec:release}) of per-pair measurements with
costs, for reproducing the contrasts, evaluating retrieval against
screen-defined labels, and studying variability.
